%% errata.tex -- supplementary material: the errata that the
%% formalization found in the fourteen sources, one by one, with the
%% Agda module in which each corrected statement is checked.
%%
%% Build: latexmk -pdf errata.tex
\documentclass[11pt]{article}
\usepackage[a4paper,margin=2.5cm]{geometry}
\usepackage{amsmath,amssymb}
\usepackage{booktabs}
\usepackage{tabularx}
\usepackage{hyperref}
\newcolumntype{Y}{>{\raggedright\arraybackslash}X}
\DeclareUnicodeCharacter{2130}{\ensuremath{\mathcal{E}}}
\DeclareUnicodeCharacter{2080}{\ensuremath{_0}}
\DeclareUnicodeCharacter{2081}{\ensuremath{_1}}
\newcommand{\ket}[1]{|#1\rangle}

\newcommand{\Z}{\mathbb{Z}}
\newcommand{\aid}[1]{\mbox{\textsf{\small #1}}}

\title{Errata Found by the Formalization\\
\large Supplementary material for ``Human-Made Scaffolds for Machine-Written
Proofs: Completeness of Quantum Circuit Equations''}
\author{Anonymous}
\date{October 2026}

\begin{document}
\sloppy
\maketitle

This document lists, source by source, every error that the
formalization found in the papers we formalized, with the corrected
statement and the Agda module in which the correction is checked. The
paper summarizes them in its Section~5 and Appendix~A; here each one is
stated in full, so that a reader can check it by hand or by a short
computation. Errors in our own transcriptions are in Section~5.2 of
the paper and are not repeated here.

Throughout, module names are relative to \path{Examples/Groups/} of
the repository, except those of the minimality development, which are
relative to \path{Examples/Independence/}.

\section{Hadamard-\texorpdfstring{$\Pi$}{Pi} (Fang, Heunen and Kaarsgaard, POPL 2026)}

References are to the extended version (arXiv:2506.06835), whose
Appendix~A contains the proof of Theorem~4.14, the completeness of
Figure~6 for $O_n(\Z[\frac1{\sqrt2}])$. Indices in this section are
the paper's, from~1.

\subsection{Lemma A.16 is false}

\paragraph{Statement.} Let $M$ be at level $(j,k,4)$ with $k\ge2$, and
let $v$ be the first four entries of column $j$, with
$\sqrt2^k v \equiv (1+\sqrt2,1,1+\sqrt2,1)$ or
$(1,1+\sqrt2,1,1+\sqrt2) \pmod 2$. The lemma says that if
\[
\sqrt2^{k-1}\,H_{[1,4]}H_{[2,3]}H_{[2,4]}H_{[2,3]}H_{[1,4]}\,v
\]
has entries in $\Z[\sqrt2]$, then so does
\[
\sqrt2^{k-1}\,H_{[1,3]}H_{[2,4]}H_{[1,2]}H_{[1,3]}H_{[2,4]}\,v .
\]

\paragraph{Counterexample.} Take $n=j=6$, $k=4$, and the unit column
\[
\sqrt2^4\cdot M_{\cdot,6} = (-1-\sqrt2,\; 1,\; -1-\sqrt2,\; 1,\; \sqrt2,\; -2+\sqrt2).
\]
It is a unit vector: the squared norm of the right-hand side is
$2(3+2\sqrt2) + 2 + 2 + (6-4\sqrt2) = 16 = (\sqrt2^4)^2$. Some matrix
of $O_6(\Z[\frac1{\sqrt2}])$ has this last column, because the
generators act transitively on unit columns. Its level is $(6,4,4)$,
and its first four entries are $\equiv(1+\sqrt2,1,1+\sqrt2,1)$ modulo
$2$, so the hypotheses hold. Applying the first word to the first four
entries gives $\sqrt2^{-3}(0,0,-1-\sqrt2,-1)$, with entries in
$\Z[\sqrt2]$ after multiplication by $\sqrt2^{3}$. Applying the second
gives
$\sqrt2^{-3}(-\tfrac{\sqrt2}{2},\,-1-\tfrac{\sqrt2}{2},\,-\tfrac{\sqrt2}{2},\,-1-\tfrac{\sqrt2}{2})$,
whose entries are not in $\Z[\sqrt2]$ after multiplication by
$\sqrt2^{3}$. So the conclusion fails.

\paragraph{Where the proof goes wrong.} The proof (page~38) writes
the entries as $a_i+\sqrt2 b_i$ and compares the $\sqrt2$-parts of the
norm equation. The constant terms of $(2a_1+1)(2b_1+1)$ and
$(2a_3+1)(2b_3+1)$ contribute $4$ to that comparison, so the bracket
the proof reduces modulo~2 is missing a $+1$. With the $+1$, the case
$a_3+a_4$ odd is no longer contradictory, and the column above lies
in it.

\subsection{Lemma A.17 is false}

Lemma~A.17 has the same hypotheses as Lemma~A.16, with $H_{[3,4]}$ in
place of $H_{[2,4]}$ in the two words, and its proof ends by appeal to
the argument of Lemma~A.16. It fails in the same way, in the case
$a_3+a_4$ even. A counterexample is the unit column
\[
\sqrt2^4\cdot M_{\cdot,6} = (1+\sqrt2,\; -1,\; 1+\sqrt2,\; 1,\; -\sqrt2,\; 2-\sqrt2),
\]
again at level $(6,4,4)$: the first word lowers the denominator
exponent of the first four entries and the second does not.

\subsection{Main Case 9 of the Main Lemma does not cover every state}

Main Case~9 of the Main Lemma (Lemma~A.7) is the case $G=H_{[1,2]}$,
$N=H_{[1,i_2]}$. In its case (2)(b)(i), second bullet, it calls
Lemma~A.19 for a path between two states that stays strictly below the
level of the edge; in the case $\ell=4$, sub-case (2), of that lemma,
items (c) and (d) rest on Lemmas~A.16 and A.17.

A state on which the proof fails is $s$ with
\[
\sqrt2^7\cdot s_{\cdot,9} = (1+3\sqrt2,\; 1,\; -3\sqrt2,\; 1+3\sqrt2,\; 2,\; 5,\; -2+\sqrt2,\; -4+2\sqrt2,\; -2-2\sqrt2),
\]
in dimension $n=9$, at level $(9,7,4)$. Algorithm~1 takes
$N=H_{[1,4]}$. The proof of Main Case~9 goes to Lemma~A.19 with
$t_2=H_{[2,6]}H_{[1,4]}\,s$ at level $(9,6,4)$ and
$(a,b,c,d)=(1,2,4,6)$. Rows 1 and 2 of $t_2$ have different residues,
which is sub-case (2) of the lemma, and only item (d) of Lemma~A.15
holds there. Main Case~9 needs a path from $t_2$ to $t_3$ strictly
below $(9,7,4)$, and every one of the eight diagrams of Lemma~A.19,
started at $t_2$ with these labels, reaches a state at level $(9,7,4)$
or $(9,8,4)$. So the proof of the Main Lemma, and with it the proof of
Theorem~4.14, does not cover this state.

The state was found by sampling random states of the relevant level
and following the case analysis of the proof on each, computing the
level of every state of every diagram. Such states are rare: none
among about 3{,}000 instances of the sub-case in dimension~7 and about
9{,}000 in dimension~8, and one among about 400 states that reach
this step in dimension~9.

\paragraph{The theorem is true.} Theorem~4.14 holds. The edge the proof
leaves open, $H$ on two odd entries of different residue classes when
exactly four entries of the pivot column are odd, is Case~3.4 of
Cl\'ement's thesis (2019), which presents the same group by other
relations. We proved his relations and those of Figure~6 equivalent
in every dimension and formalized his Case~3.4 for this edge; on the
state $s$ above his diagram~(32), on rows $1,2,4,6$ with the further
rows $3$ and $5$, gives a path that never rises above the level of
$s$. Modules: \path{Real-Clifford+CH-TwoLevel/Presentation} states
the theorem with the hard edge as a hypothesis
(\aid{MainLemma.Hard}), and
\path{Real-Clifford+CH-TwoLevel-Clement/Presentation} discharges it
(\aid{hard}) and concludes \aid{presentation} unconditionally;
\path{Real-Clifford+CH-TwoLevel/Equivalence} is the equivalence of
the two relation sets.

\subsection{Citations in the proof of Lemma A.19(d)}

The proof of item (d) of Lemma~A.19 cites Propositions~A.28 and A.29.
Proposition~A.29 as printed repeats Proposition~A.22. What the step
needs is the relation of Proposition~A.28 and a second relation that
is stated nowhere in the paper. That relation holds, and it follows
from Proposition~A.26 in the same way as A.28 follows from A.27. (In
the formalization this part of Lemma~A.19 is not used, the case being
handled by Cl\'ement's argument.)

\section{Real Clifford+\texorpdfstring{$CH$}{CH} circuits (Cl\'ement, 2026)}

References are to arXiv:2602.06644.

\paragraph{Dependency.} Theorem~4.4 of the paper is Theorem~4.14 of
Hadamard-$\Pi$, whose proof has the gap above. The paper's main
theorem, Theorem~8.9, depends on it through Theorem~4.10 as well. Both
are now discharged: \path{Real-Clifford+CH/Presentation} proves
Theorem~8.9 with no hypotheses, with Theorem~4.4 taken from the
Cl\'ement development (\path{Auxiliary/Theorem44Proof}).

\paragraph{Equations (36) and (37).} These write a Hadamard pair with
an $H$-pattern through a network of swaps and negations. As printed,
the negation of a control bit $x$ is $\mathcal{E}^{1-x}$, an $X$ gate
where the bit is $0$. The equations hold with $\mathcal{E}^{x}$, an
$X$ gate where the bit is $1$, and fail at every instance the other
way round. Checked by deciding both equations on three qubits at every
admissible choice of indices (\path{Real-Clifford+CH/Auxiliary/Checks8});
the transcription with $\mathcal{E}^{x}$ is \aid{E-X} and \aid{ℰ} of
\path{Real-Clifford+CH/Encoding}.

\paragraph{Side conditions of Figure 8.} The rules (30), (34) and (42)
of Figure~8 need their $X$ and $H$ letters to have distinct indices, a
convention the paper applies throughout but does not state in the
figure. Without the conditions the rules are unsound, which the same
decision procedure shows. They are stated as hypotheses in
\path{Real-Clifford+CH/Auxiliary/Figure8}.

\section{Real stabilizer circuits (Makary, Ross and Selinger, QPL 2021)}

References are to arXiv:2109.05655 and its supplement of derivations.
Below, a relation between circuits is written as an equation of words
over the gates $H_i$, $Z_i$ and $CZ_{i,i+1}$ on wires numbered from
the bottom, in operator order: the rightmost gate acts first, as in a
matrix product. These are the relations as our normalization uses
them (\path{Real-Clifford/Relations/Rules}), and each is checked by
evaluation on matrices.

\subsection{Eleven typed relations of Appendix B}

Of the 137 typed relations of Appendix~B, eleven are false as printed:
the two sides have different matrices, not only up to a sign. Six have
a wrong right-hand side, and five lack a Pauli gate at the end of the
right-hand side. Table~\ref{tab:typed} lists them, and the corrected
relations follow as word equations; $t$ and $b$ denote the top and
bottom wire of a two-wire gate in the paper's notation.

\begin{table}[h]
\caption{The eleven misprinted typed relations.}
\label{tab:typed}
\begin{tabularx}{\linewidth}{@{}l l Y@{}}
\toprule
Relation & Error & Note \\
\midrule
$H;A_3$ & right side & printed $A_2;H$; correct is $A_3;H$ \\
$X;C_1$ & right side & printed $C_1;Z$; correct is $C_2$ \\
$X;C_2$ & right side & \\
$H_0;B_6$, $H_0;B_7$, $H_0;B_8$ & right side & \\
$CZ;A_3;B_7$ & Pauli tail & \\
$Z_t;B_2$ & Pauli tail & printed $B_2;X_t$; correct is $B_2;X_t;Z_b$ \\
$Z_t;B_6$ & Pauli tail & \\
$SXZC;B_5$, $SXZC;B_8$ & Pauli tail & \\
\bottomrule
\end{tabularx}
\end{table}

\begin{align*}
H;A_3:\quad & H_0 = H_0 \\
X;C_1:\quad & H_0Z_0H_0 = H_0Z_0H_0 \\
X;C_2:\quad & H_0Z_0H_0\,H_0Z_0H_0 = \varepsilon \\
H_0;B_6:\quad & CZ_{01}H_1H_0CZ_{01}H_0 = H_1H_0CZ_{01}H_1CZ_{01}H_1H_0CZ_{01}H_1H_0CZ_{01} \\
H_0;B_7:\quad & CZ_{01}H_1H_0CZ_{01}H_1H_0 = H_1H_0CZ_{01}H_1CZ_{01}H_1H_0CZ_{01}H_1H_0CZ_{01}H_1 \\
H_0;B_8:\quad & CZ_{01}H_1H_0CZ_{01}H_0CZ_{01}H_0H_0 \\
 & \quad = CZ_{01}H_1Z_1H_1H_0CZ_{01}CZ_{01}H_1H_0CZ_{01}H_0CZ_{01}H_0 \\
CZ;A_3;B_7:\quad & CZ_{01}H_1H_0CZ_{01}H_1CZ_{01} \\
 & \quad = H_0Z_0H_0\,H_0CZ_{01}H_1CZ_{01}H_1H_0CZ_{01}H_1H_0CZ_{01}H_1H_0CZ_{01}H_0 \\
Z_t;B_2:\quad & CZ_{01}H_1H_0CZ_{01}Z_1 = Z_0H_1Z_1H_1CZ_{01}H_1H_0CZ_{01} \\
Z_t;B_6:\quad & CZ_{01}H_1H_0CZ_{01}Z_1 = Z_0H_1Z_1H_1CZ_{01}H_1H_0CZ_{01} \\
SXZC;B_5:\quad & CZ_{12}H_2H_1CZ_{12}H_2H_1CZ_{12}H_1CZ_{01}H_1CZ_{01}H_1 \\
 & \quad = Z_0Z_0H_0CZ_{01}H_0H_1CZ_{01}H_1CZ_{12}H_2CZ_{12}H_2H_1 \\
 & \qquad\quad CZ_{01}H_1H_0CZ_{01}H_0CZ_{12}H_2H_1CZ_{12}H_2H_1CZ_{12}H_1 \\
SXZC;B_8:\quad & CZ_{12}H_2H_1CZ_{12}H_1CZ_{12}H_1CZ_{01}H_1CZ_{01}H_1 \\
 & \quad = Z_0CZ_{01}H_1CZ_{01}H_1H_0CZ_{01}H_0H_1CZ_{01}H_1CZ_{12}H_1 \\
 & \qquad\quad CZ_{01}H_1H_0CZ_{01}H_0CZ_{12}H_2H_1CZ_{12}H_1CZ_{12}H_1
\end{align*}

The words are those of \aid{Rule-HA-3}, \aid{Rule-XC-1},
\aid{Rule-XC-2}, \aid{Rule-H0B-6}, \aid{Rule-H0B-7}, \aid{Rule-H0B-8},
\aid{Rule-CZAB-3-7}, \aid{Rule-Z1B-2}, \aid{Rule-Z1B-6},
\aid{Rule-SXZCB-5} and \aid{Rule-SXZCB-8} in
\path{Real-Clifford/Relations/Rules}; the typed forms, with the
layer types, are the corresponding entries of
\path{Real-Clifford/Tables}, each paired with its relation so that a
wrong entry does not typecheck. The normal form, its uniqueness and
the count of Corollary~4.16 go through with the corrected relations.

\subsection{Two smaller problems in the paper}

\begin{itemize}
\item Figure~10 prints one equation without its ``$=$''.
\item Definition~5.3 of the normal forms does not allow a $CZ$ whose
  top wire is a single wire labelled~2, although the left-hand sides
  of rules 5.1--5.4 have exactly such a gate. The inductive normal
  forms of \path{Real-Clifford/NormalForm} include the case.
\end{itemize}

\subsection{The supplement}

Six derivations in the supplement reach one of the eleven false
conclusions above; the formalization derives the corrected relations
instead (\path{Real-Clifford/Relations/Typed} and
\path{Derived}). Besides these, the supplement has five misprinted
intermediate rows, four steps labelled with the wrong axiom, and some
forty citations of typed relations by an older figure numbering; the
script that followed the derivations line by line stopped at each of
these, and continued with the step that does hold. The script was not
kept, and we did not record the page of each of these smaller slips;
the counts are those reported when the modules were generated.

\section{Clifford circuits (Selinger, LMCS 2015)}

References are to arXiv:1310.6813 and its 105-page supplement, whose
derivations are followed line by line in
\path{Qubit-Clifford/Relations}. The paper's figures all check.
Three citations in the supplement are wrong, each in a step whose
conclusion is right:

\begin{itemize}
\item The proof of Lemma~1.3 says that property (19) is relation
  (C10); it is (C12).
\item Two steps in the proof of equation (21) are labelled (11). Rule
  (11) cannot reorder two $CZ$ gates sharing a wire; the steps use
  (19).
\end{itemize}

These are the only wrong citations among the derivations of
Lemmas~1.1--1.3 and equations (3.1)--(3.101).

\section{\texorpdfstring{$O_n(\Z[\frac12])$}{On(Z[1/2])} and \texorpdfstring{$L_n$}{Ln} (Li, Ross and Selinger, QPL 2021)}

References are to arXiv:2106.01175. Relation (7d) of Table~2, which
conjugates $X_{[a,a+1]}$ by $I\otimes H$, has its two right-hand sides
exchanged. With indices from~1, as in the paper: for $a$ odd the block
$\{a,a+1\}$ is one block of $I\otimes H$, and
$(I\otimes H)\,X_{[a,a+1]}\,(I\otimes H) = (-1)_{[a+1]}$; for $a$ even
it is $X_{[a,a+1]}\,K_{[a-1,a,a+1,a+2]}$. The printed right-hand side,
$(-1)_{[a+1]}^{\,a+1}\,X_{[a,a+1]}^{\,a}\,K_{[a-1,a,a+1,a+2]}^{\,a}$ with
exponents read modulo~2, has the two cases the other way round. Found
when the relation failed to check numerically; the corrected forms are
\aid{r7d₀} and \aid{r7d₁} (indices from~0) in
\path{CCX+HH-TwoLevel/Scaled/Syntactics}, and both are proved sound
in \path{CCX+HH-TwoLevel/Scaled/Table2}.

\section{Simpler presentations (Blake, FSCD 2026)}

References are to arXiv:2602.09874.

\paragraph{The scalar of rule (E).} Figure~4 gives the scalar of the
qutrit Clifford rule (E) as $\omega_{12}^7$. With the semantics of
Section~3.2 ($H\ket{x} = \sum_k e^{2\pi i xk/3}\ket{k}/\sqrt3$,
$S = \mathrm{diag}(1,1,e^{2\pi i/3})$), $SHS$ equals
$\omega_{12}\,H^3S^2H^3$, and $\omega_{12}^7 = -\omega_{12}$ makes the
rule unsound. It is sound with $\omega_{12}$, which is what
\path{Qutrit-Clifford/Syntactics} states and
\path{Qutrit-Clifford/Semantics} checks. No separating
interpretation of Figure~8(c) sees the exponent, so the minimality
theorem is unaffected.

\paragraph{A separator that is not a model.} The separating
interpretation $\#\{CX,\sigma\}[2]$ used for the independence of the
Clifford rule (CZ) does not satisfy the three-wire rule (I), whose
sides carry three $CNOT$s and two swaps against two $CNOT$s
(\aid{Sep-CZ.not-a-model-of-I} in \path{Clifford/Minimality}). The
independence claim survives: a derivation between two-wire circuits
uses no three-wire rule, the paper's Remark~14, which is proved as
\aid{Prop.truncate} and applied to every bounded separator.

\end{document}
